\section{Conclusiones}

De forma experimental, se determinó que la densidad del agua mediante el principio de Arquímedes fue $\rho_f=(0{,}991\pm0{,}078)\,\text{g/cm}^3$. Este resultado presentó un error porcentual de $0{,}86\%$ respecto al valor de referencia de $1{,}00\,\text{g/cm}^3$ y un coeficiente de determinación de $R^2=0{,}9930$, lo que indicó una buena correlación entre la diferencia de masa y la longitud de varilla sumergida. Por tanto, se cumplió el objetivo de calcular la densidad del fluido a partir del empuje.

En la segunda parte, los tiempos de vaciado presentaron errores entre $5{,}73\%$ y $22{,}48\%$ respecto a los valores teóricos obtenidos mediante el teorema de Torricelli. Aunque estos errores evidenciaron desviaciones respecto al modelo ideal, se observó que el tiempo de vaciado disminuyó al reducirse la altura inicial del agua, de acuerdo con la tendencia esperada.

En conclusión, se cumplieron los objetivos principales asociados con ambos experimentos. El principio de Arquímedes se corroboró cuantitativamente dentro de la incertidumbre experimental, mientras que el teorema de Torricelli reprodujo la tendencia esperada, aunque con discrepancias atribuibles a condiciones no ideales del montaje.

Como mejora para trabajos futuros, se recomienda utilizar una balanza digital para aumentar la precisión de las mediciones y un cronómetro de mayor resolución.
